\documentclass[UTF8, 12pt, a4paper]{ctexart}

% ==================== 宏包加载 ====================
\usepackage{geometry}
\usepackage{amsmath}
\usepackage{graphicx}
\usepackage{booktabs}
\usepackage{tabularx}
\usepackage{hyperref}
\usepackage{xcolor}
\usepackage{caption}
\usepackage{cite}
\usepackage{makecell}
\usepackage{multirow}
\usepackage{enumitem}
\usepackage{array}
\usepackage{float}
\usepackage{titlesec}

% 【核心新增神器】：tcolorbox 用于绘制支持自动跨页的精美边框
\usepackage[most]{tcolorbox}

% ==================== 全局格式设置 ====================
\geometry{left=2.5cm, right=2.5cm, top=2.8cm, bottom=2.8cm}

% 定义主题色
\definecolor{themeblue}{RGB}{0, 82, 155}
\definecolor{themegray}{RGB}{90, 90, 90}
\definecolor{lightgray}{RGB}{245, 245, 245}
\definecolor{tablerule}{RGB}{60, 60, 60}
\definecolor{sectioncolor}{RGB}{0, 70, 140}

\hypersetup{
    colorlinks=true,
    linkcolor=themeblue,
    citecolor=themeblue,
    urlcolor=themeblue
}

\captionsetup{
    font={small, sf},
    labelfont={bf, color=themeblue},
    labelsep=quad,
    skip=8pt
}
\graphicspath{{images/}}

\newcommand{\mybody}{\songti\fontsize{10.5bp}{18bp}\selectfont\color{black}}
\newcolumntype{C}{>{\centering\arraybackslash}X}

% ==================== 页眉页脚设置 ====================
\usepackage{fancyhdr}
\pagestyle{fancy}
\fancyhf{}
\renewcommand{\headrulewidth}{0.6pt}
\renewcommand{\footrulewidth}{0pt}
\fancyhead[C]{\small\color{themegray} 宇宙无敌暴龙战士队 \quad|\quad 王睿恒（队长）\quad 杨林涛 \quad 王彬玮}
\fancyfoot[C]{\small\color{themegray} \thepage}

% ==================== 1. 带表头的标准报告外框 ====================
\newtcolorbox{reportframe}[1][]{
    enhanced,
    breakable,
    colback=white,
    colframe=black,
    boxrule=0.6pt,
    arc=0pt,
    outer arc=0pt,
    boxsep=0pt,
    left=1.8em, right=1.8em, top=1.5em, bottom=1.5em,
    title={%
        \renewcommand{\arraystretch}{1.8}%
        \setlength{\arrayrulewidth}{0.6pt}%
        \begin{tabularx}{\linewidth}{@{} >{\hsize=1.4\hsize\centering\arraybackslash}X | >{\hsize=0.6\hsize\centering\arraybackslash}X @{}}
            \multicolumn{2}{c}{\zihao{-4} \bfseries 第 21 届 "中控杯" 浙江大学机器人竞赛} \\
            \hline
            \multirow{2}{*}{\zihao{1} \bfseries 技术方案文档} & {\zihao{5} 参赛项目} \\
            \cline{2-2}
            & {\zihao{5} 对抗赛道}
        \end{tabularx}%
    },
    coltitle=black,
    colbacktitle=white,
    titlerule=0.6pt,
    toptitle=0pt,
    bottomtitle=0pt,
    fonttitle=\normalfont,
    #1
}

% ==================== 2. 无表头的续页纯净外框 ====================
\newtcolorbox{reportframe_notitle}[1][]{
    enhanced,
    breakable,
    colback=white,
    colframe=black,
    boxrule=0.6pt,
    arc=0pt,
    outer arc=0pt,
    boxsep=0pt,
    left=1.8em, right=1.8em, top=1.5em, bottom=1.5em,
    #1
}

% ==================== 自定义节标题命令 ====================
\newcommand{\sectiontitle}[1]{%
    \vspace{0.5em}%
    {\zihao{4}\bfseries\color{sectioncolor} #1}%
    \vspace{0.3em}%
    \par\noindent\textcolor{sectioncolor}{\rule{\linewidth}{0.8pt}}%
    \vspace{0.8em}%
}

% ==================== 自定义小节标题命令 ====================
\newcommand{\subsectiontitle}[1]{%
    \vspace{0.6em}%
    \noindent{\color{sectioncolor}\vrule width 3pt height 1em depth 0.2em\hspace{0.5em}}%
    \textbf{#1}%
    \vspace{0.3em}\par%
}

% ==================== 正文部分 ====================
\begin{document}

% -------------------- 封面 --------------------
\begin{figure}[htbp]
    \centering
    \includegraphics[width=0.8\textwidth]{image1.png}
\end{figure}

\begin{center}
    \vspace*{0.5cm}
    {\zihao{1} \bfseries 第 21 届} \\[1.2cm]
    {\zihao{1} \bfseries "中控杯" 浙江大学机器人竞赛} \\[1.5cm]
    {\zihao{2} 创新机器人制作竞赛} \\[0.6cm]
    {\zihao{2} 创新组技术方案设计报告}
    \vspace{1cm}
\end{center}

\begin{center}
    \renewcommand{\arraystretch}{2.0}
    \setlength{\arrayrulewidth}{0.5pt}
    \begin{tabularx}{0.85\textwidth}{>{\centering\arraybackslash}p{4cm} C}
        \toprule[1.2pt]
        {\zihao{3} 团队名称} & {\zihao{3} 宇宙无敌暴龙战士队} \\
        \midrule
        {\zihao{3} 团队编号} & {\zihao{3} 8182} \\
        \midrule
        {\zihao{3} 队长姓名} & {\zihao{3} 王睿恒} \\
        \midrule
        {\zihao{3} 队员 1 姓名} & {\zihao{3} 杨林涛} \\
        \midrule
        {\zihao{3} 队员 2 姓名} & {\zihao{3} 王彬玮} \\
        \midrule
        {\zihao{3} 队员 3 姓名} & {\zihao{3} 无} \\
        \midrule
        {\zihao{3} 指导教师} & {\zihao{3} 无} \\
        \bottomrule[1.2pt]
    \end{tabularx}
\end{center}

\clearpage

% -------------------- 机械设计 --------------------
\begin{reportframe}
    \sectiontitle{1、机器人机械设计方案}

    \mybody
    本机器人采用子母车架构，面向对抗赛"争分夺秒抢夺物资"的核心需求，将机械系统拆分为三大功能模块并行协作：\textbf{大车（母车）}搭载四组麦克纳姆轮与推块臂，负责批量推送方块至高分区，同时利用体型优势封锁对手路径；\textbf{小车（子车）}可脱离母车独立行动，深入场地各区域补充推块得分；\textbf{机械臂}安装于大车上方，专职抓取圆环套入 2 分区立柱。三者在比赛中按阶段分离、并行得分，最大化单位时间得分效率。整体采用双层 6061-T6 铝合金板材框架，上下层通过 M3 铜柱螺接，子车经顶板腰型槽与电控锁紧机构固定于母车之上，比赛中由主控信号触发释放。

    \begin{center}
        \vspace{0.5em}
        \begin{minipage}{0.48\textwidth}
            \centering
            \includegraphics[height=0.22\textheight]{image3.png}
            \vspace{0.5em}
            \captionof{figure}{下车底板结构示意图}
        \end{minipage}
        \hfill
        \begin{minipage}{0.48\textwidth}
            \centering
            \includegraphics[height=0.22\textheight]{image2.png}
            \vspace{0.5em}
            \captionof{figure}{下车顶板结构示意图}
        \end{minipage}
        \vspace{0.5em}
    \end{center}

    \subsectiontitle{（1）底板设计（见图 1）}
    底板作为下车的动力基准，采用 6061-T6 铝合金板材激光切割而成。其表面预留了高密度的安装孔位阵列，四组驱动电机座通过 M3 螺栓固定于底板预留孔位，万向轮支撑件经定位销与螺栓双重定位连接。蓄电池通过扎带与防滑垫固定于底板中央低位区域，使整车重心处于安全包络内，显著提升了高速转向时的动态稳定性。底板前端预留了推块臂的安装基准面，使推块臂与底盘驱动形成刚性整体，保证批量推送方块时的力传递效率。

    \subsectiontitle{（2）顶板设计（见图 2）}
    顶板兼具结构加强与子车承载功能，通过四根 M3 铜柱与底板螺接形成刚性双层框架。子车通过腰型槽配合电磁锁紧机构固定于顶板上表面，比赛中由主控 GPIO 信号触发释放，小车依靠自身动力驶离母车展开独立作业。
    \begin{itemize}[leftmargin=2em, itemsep=3pt, parsep=0pt, topsep=4pt]
        \item \textbf{功能腰型槽}：中心布置的三组腰型槽在实现拓扑轻量化的同时，为子车的锁紧机构提供了 $\pm 25$\,mm 的机械调节冗余，增强了系统适配性。
        \item \textbf{90$^\circ$ 折弯工艺}：侧边采用折弯补强结构，极大地提高了板材的抗弯惯性矩，有效防止子车在载入对接过程中因局部受压而产生形变。
    \end{itemize}

    \subsectiontitle{（3）底盘工艺指标}
    板材选用 6061-T6 铝合金，加工公差控制在 $\pm 0.1$\,mm 级别。所有边缘均经过倒角处理，并预留了电气走线孔，实现了机械结构与电路系统的模块化集成。

    \subsectiontitle{（4）专用物块推动机构设计（见图 3）}
    推动机构（推块臂）安装于大车前端，是实现\textbf{批量推块得分}的核心执行部件。在对抗赛中，推块臂可一次接触并推送 2—3 个方块，将其从场地中央直线推入 3 分区（3 分/个），得分效率远高于逐个抓取。为保证多块同推时的一致性，整体采用了非线性几何与严苛的公差设计。

    \begin{center}
        \vspace{0.5em}
        \includegraphics[height=0.22\textheight]{image4.png}
        \vspace{0.5em}
        \captionof{figure}{物块推动机构三维结构示意图}
        \vspace{0.3em}
    \end{center}

    \begin{itemize}[leftmargin=2em, itemsep=3pt, parsep=0pt, topsep=4pt]
        \item \textbf{多节弯曲避障臂}：臂体采用特定的非线性弯曲几何形状，在有效避让物块上方潜在障碍物的同时，优化了刚性分布。经有限元分析（FEA）验证，在受侧向推力时其弹性形变 $\delta \leq 0.05$\,mm。
        \item \textbf{高精度法兰与执行推块}：顶部法兰通过四个 M4 沉头螺栓与车体刚性连接，保证极大的扭转惯性矩。底部执行推块（基准面 1）的宽度公差严格控制在 $\pm 0.50$\,mm 以内，接触面垂直度公差 $\leq \pm 0.1^\circ$，确保动作的正交性与一致性。
        \item \textbf{材料与表面处理}：主体采用 6061-T6 铝合金进行五轴 CNC 精密加工，以降低运动惯量。接触面经阳极氧化处理并添加 PTFE（聚四氟乙烯）涂层，显著降低磨损，保障长期作业精度。
    \end{itemize}

    \vspace*{\fill}
\end{reportframe}

\clearpage

% -------------------- 电路设计 --------------------
\begin{reportframe}[breakable=false]
    \sectiontitle{2、机器人电路设计方案}

    \mybody
    电路系统是机器人的"神经与血管"，负责能源供给、信号处理与运动控制。整体采用\textbf{"电源$\,\to\,$主控$\,\to\,$驱动$\,\to\,$感知"}四层架构，各层之间通过标准化接口连接，便于独立调试与模块化更换。

    \subsectiontitle{（1）电源系统}
    采用格式锂电池作为全车唯一动力源，电池通过 XT60 接头连接至电源分配板，分为两条独立回路：\textbf{动力回路}直接为 AT8236 电机驱动模块供电；\textbf{逻辑回路}经 AT8236 板载稳压电路降压至 5\,V / 3.3\,V，分别为舵机与 STM32 主控、K230 视觉模块供电。蜂鸣器电量显示器并联于电池端，实时监测电压，低压时自动报警，防止比赛中意外断电。

    \subsectiontitle{（2）主控系统}
    以 \textbf{STM32H743ZIT6} 为主控核心，运行巡线 PID、状态机调度与视觉数据解析等全部上层逻辑；\textbf{STM32H750VBT6} 作为副控，专职负责机械臂步进电机脉冲输出与舵机 PWM 生成。主控与副控之间通过 UART 串口（TTL 电平直连）交换状态指令与坐标数据。\textbf{K230 视觉模块}同样经 UART 串口接入主控，以固定帧率输出目标物体的像素坐标。此外，主控通过一路 GPIO 控制顶板电磁锁紧机构的通断，实现对子车的定时释放。

    \subsectiontitle{（3）驱动系统}
    底盘四轮各配一台 520 直流减速电机（内置霍尔编码器），由两块 \textbf{AT8236 双路驱动模块}分别驱动。编码器 A/B 相信号接入 STM32 定时器编码器接口，实现四轮独立闭环调速。机械臂夹取由\textbf{数字舵机}完成，通过主控 PWM 输出精确控制张合角度；升降机构采用步进电机驱动丝杆，由副控脉冲方向接口控制行程。

    \subsectiontitle{（4）通信与安全}
    \textbf{315\,MHz 四路遥控模块}用于远程启停与紧急制动，接收端 IO 直连主控 GPIO 中断引脚，响应延迟 $<\!10$\,ms。整车线缆采用端子化接插件（XH2.54 / PH2.0），便于快速拆装与故障排查。子母车分离后，双方按预设时间表独立执行各自任务路径，无需实时通信，降低了系统复杂度。

    \vspace{0.5em}
    以下为核心电子元器件实物清单：

    \vspace*{\fill}
\end{reportframe}
\clearpage

% -------------------- 电路设计（元器件图片页1） --------------------
\begin{reportframe_notitle}[breakable=false]

    % ==================== 第一排 ====================
    \noindent
    \begin{minipage}[t]{0.48\textwidth}
        \centering
        \includegraphics[height=6.2cm, keepaspectratio]{images/315mHz四路.png}\\[0.2em]
        \captionof{figure}{315\,mHz 四路}
    \end{minipage}%
    \hfill
    \begin{minipage}[t]{0.48\textwidth}
        \centering
        \includegraphics[height=6.2cm, keepaspectratio]{images/520霍尔编码器直流减速机.png}\\[0.2em]
        \captionof{figure}{520 霍尔编码器直流减速机}
    \end{minipage}

    \vspace{0.4em}

    % ==================== 第二排 ====================
    \noindent
    \begin{minipage}[t]{0.48\textwidth}
        \centering
        \includegraphics[height=6.2cm, keepaspectratio]{images/蜂鸣器1-8s二合一电量显示器.png}\\[0.2em]
        \captionof{figure}{蜂鸣器 1-8s 二合一电量显示器}
    \end{minipage}%
    \hfill
    \begin{minipage}[t]{0.48\textwidth}
        \centering
        \includegraphics[height=6.2cm, keepaspectratio]{images/格式电池.png}\\[0.2em]
        \captionof{figure}{格式电池}
    \end{minipage}

    \vspace{0.4em}

    % ==================== 第三排 ====================
    \noindent
    \begin{minipage}[t]{0.48\textwidth}
        \centering
        \includegraphics[height=6.2cm, keepaspectratio]{images/六轴机械臂抓手.png}\\[0.2em]
        \captionof{figure}{六轴机械臂抓手}
    \end{minipage}%
    \hfill
    \begin{minipage}[t]{0.48\textwidth}
        \centering
        \includegraphics[height=6.2cm, keepaspectratio]{images/数字舵机.png}\\[0.2em]
        \captionof{figure}{数字舵机}
    \end{minipage}

    \vspace*{\fill}
\end{reportframe_notitle}
\clearpage

% -------------------- 电路设计（元器件图片页2） --------------------
\begin{reportframe_notitle}[breakable=false]

    % ==================== 第一排 ====================
    \noindent
    \begin{minipage}[t]{0.48\textwidth}
        \centering
        \includegraphics[height=6.2cm, keepaspectratio]{images/直流减速电机GMR霍尔编码器.png}\\[0.2em]
        \captionof{figure}{直流减速电机 GMR 霍尔编码器}
    \end{minipage}%
    \hfill
    \begin{minipage}[t]{0.48\textwidth}
        \centering
        \includegraphics[height=6.2cm, keepaspectratio]{images/AT8236双路驱动模块带稳压板.png}\\[0.2em]
        \captionof{figure}{AT8236 双路驱动模块带稳压板}
    \end{minipage}

    \vspace{0.4em}

    % ==================== 第二排 ====================
    \noindent
    \begin{minipage}[t]{0.48\textwidth}
        \centering
        \includegraphics[height=6.2cm, keepaspectratio]{images/K230视觉模块.png}\\[0.2em]
        \captionof{figure}{K230 视觉模块}
    \end{minipage}%
    \hfill
    \begin{minipage}[t]{0.48\textwidth}
        \centering
        \includegraphics[height=6.2cm, keepaspectratio]{images/STM32H743ZIT6核心板.png}\\[0.2em]
        \captionof{figure}{STM32H743ZIT6 核心板}
    \end{minipage}

    \vspace{0.4em}

    % ==================== 第三排 ====================
    \noindent
    \begin{minipage}[t]{0.48\textwidth}
        \centering
        \includegraphics[height=6.2cm, keepaspectratio]{images/金属齿轮.png}\\[0.2em]
        \captionof{figure}{25KG 金属齿轮数字舵机}
    \end{minipage}%
    \hfill
    \begin{minipage}[t]{0.48\textwidth}
        \centering
        \includegraphics[height=6.2cm, keepaspectratio]{images/STM32H750VBT6开发板.png}\\[0.2em]
        \captionof{figure}{STM32H750VBT6 开发板}
    \end{minipage}

    \vspace*{\fill}
\end{reportframe_notitle}

\clearpage

% -------------------- 第四部分：软件设计 --------------------
\begin{reportframe}
    \sectiontitle{3、机器人软件设计方案}

    \mybody
    软件系统是整套子母车协同策略的"大脑"，在有限的比赛时间内调度大车推块、小车独立得分与机械臂圆环操作三条并行任务线。主要实现五大核心功能：\textbf{全向底盘运动控制}、\textbf{红外巡线导航}、\textbf{视觉目标识别与定位}、\textbf{机械臂抓取控制}以及\textbf{子车释放与并行调度}。系统基于 STM32H743ZIT6 平台，采用 C 语言在 Keil MDK 环境下开发，以模块化方式组织代码，由主程序状态机统一调度。设计目标为实现全自主完赛，巡线偏差控制在 5\,mm 以内，视觉定位精度优于 1\,cm。

    \subsectiontitle{（1）整体架构}
    主程序运行于 STM32H743ZIT6 主控芯片上，通过有限状态机管理机器人行为。针对对抗赛策略需求，定义了完整的运行状态集：\texttt{SPRINT\_PUSH}（闪电推块）、\texttt{DEPLOY\_CHILD}（释放子车）、\texttt{PUSH\_BATCH}（批量推块）、\texttt{RING\_HUNT}（圆环搜索）、\texttt{RING\_PLACE}（圆环套杆）、\texttt{DEFEND}（防守封锁）及 \texttt{OPPONENT\_AVOID}（对手规避）等。每个控制周期（10\,ms）根据当前状态与传感器输入调用对应功能函数。副控 STM32H750VBT6 接收主控指令，专职执行机械臂运动控制。

    \subsectiontitle{（2）底盘运动控制}
    基于麦克纳姆轮运动学模型，推导四轮转速与整车速度向量（$v_x$、$v_y$、$\omega$）之间的映射关系，封装为 \texttt{Chassis\_SetVelocity()} 函数。结合 MG513 编码器的转速反馈，对各轮分别实施 PID 闭环调速，保证直行与全向移动的精度。

    \subsectiontitle{（3）巡线模块}
    读取底盘前端 8 路红外循迹传感器的数字电平，通过加权偏差算法计算当前偏移量，将偏差输入 PID 控制器，输出左右侧电机的差速修正量，实现稳定巡线。当传感器检测到十字路口特征（多路同时触发）时，触发路口判断逻辑，执行预设的直行、转弯或停车动作。

    \subsectiontitle{（4）视觉融合}
    K230 视觉模块通过串口中断方式向 STM32 发送目标坐标帧。主控解析坐标后，在视觉追踪状态下计算目标与图像中心的像素偏差，换算为底盘横移与旋转指令，驱动小车对准目标，随后切换至机械臂抓取状态。此外，K230 还承担对手机器人的实时检测任务：当视野中出现对手时，主控自动切换至 \texttt{OPPONENT\_AVOID} 状态，麦克纳姆轮横移避让或主动封锁，确保己方得分路径畅通。

    \subsectiontitle{（5）机械臂控制}
    升降步进电机通过脉冲数控制丝杆位移，封装为 \texttt{Arm\_MoveTo(height)} 函数；舵机夹取通过 \texttt{Servo\_SetAngle()} 设定张合角度。抓取流程为：下降至目标高度 $\rightarrow$ 舵机夹紧 $\rightarrow$ 上升至安全高度 $\rightarrow$ 运送至目标区域 $\rightarrow$ 舵机释放。

    \subsectiontitle{（6）子车释放与独立调度}
    大车到达预设分离点后，主控触发电磁锁释放子车。子车搭载独立的巡线传感器与驱动电机，脱离后按预设路径自主巡线，执行补充推块或区域占位任务。子母车采用"时间表盲协同"策略——双方在分离前即确定各自任务路径与时间节点，无需分离后实时通信，降低系统复杂度并提升可靠性。

    \subsectiontitle{（7）已完成功能}
    目前已完成巡线、路口判断、底盘全向运动、视觉串口解析等基础功能，机械臂抓取流程、子车释放调度与完整状态机仍在调试优化中，以最终参赛代码为准。

    \vspace*{\fill}
\end{reportframe}

\clearpage

% -------------------- 第五部分：流程与实现 --------------------
\begin{reportframe}
    \sectiontitle{4、机器人完赛流程及基本功能实现方案}

    \mybody
    本机器人采用子母车架构，针对对抗赛"红蓝双方同场争夺物资"的特点，制定了\textbf{大车批量推块 + 小车独立得分 + 机械臂圆环操作}三线并行的高效完赛策略。大车凭借推块臂的批量推送能力主攻 3 分区高分方块，小车脱离后独立补充推块并实施区域占位，机械臂专职完成圆环套杆操作。整体流程分为四个阶段：

    \subsectiontitle{（1）闪电开局阶段（0—40 秒）}
    315\,MHz 遥控触发启动后，STM32 完成系统自检（$<\!1$\,s）。大车全速沿预设路径巡线前进至方块集中区域，推块臂对准方块群，利用麦克纳姆轮横向微调对齐后，一次推送 2—3 个方块直线推入最近的 3 分区。此阶段小车固定于母车上随行，机械臂在推块行进路线上"顺路"抓取第一个己方圆环，不专门绕路。

    \subsectiontitle{（2）分兵作战阶段（40 秒—1 分 40 秒）}
    大车到达预设分离点后，主控触发电磁锁释放小车。此后三条得分线并行运作：\textbf{大车}继续执行批量推块任务，将剩余方块推入 3 分区或 1 分区，同时利用体型优势对对手实施合法路径封锁；\textbf{小车}脱离后独立巡线，深入场地各区域补充推块得分，或前往对手半场抢夺"无主方块"；\textbf{机械臂}利用大车推块到位后的返程间隙，在 2 分区立柱完成圆环套杆操作。

    \subsectiontitle{（3）高效收割阶段（1 分 40 秒—2 分 40 秒）}
    大车通过 K230 视觉广域扫描场地剩余方块分布，将零散方块批量推入最近的得分区（优先填满 3 分区剩余空位，其次推入 1 分区）。机械臂完成剩余己方圆环的套杆操作。小车继续独立推块，若己方已取得明显优势则转入防守模式，占据对手前往得分区的关键路口。

    \subsectiontitle{（4）收尾锁定阶段（2 分 40 秒—比赛结束）}
    大车与小车协同执行收尾策略：领先时，两车分别占据关键得分区入口附近，合法阻碍对手放置；落后时，全力冲刺将任何可触及的方块推入最近的 1 分区，争取每一分。

    \subsectiontitle{（5）核心功能实现要点}
    \begin{itemize}[leftmargin=2em, itemsep=3pt, parsep=0pt, topsep=4pt]
        \item \textbf{批量推块效率}：推块臂宽度覆盖 2—3 个方块，配合麦克纳姆轮横向对齐，单次推送得分可达 6—9 分（3 分区），效率远超逐个抓取。
        \item \textbf{巡线稳定性}：采用 PID 差速控制与 8 路红外传感器阵列，确保在不同路面条件下可靠追线。
        \item \textbf{视觉多目标识别}：K230 模块同时检测方块、圆环与对手机器人，主控根据识别结果动态切换推块/抓取/避让状态。
        \item \textbf{机械臂圆环操作}：步进电机配合限位检测保证套杆高度一致；舵机张合角度精确匹配圆环内径，套杆成功率 $>\!90\%$。
        \item \textbf{子母车并行调度}：分离后双车按预设时间表独立执行任务，互不依赖实时通信，实现得分通道倍增。
    \end{itemize}

    \vspace{0.8em}
    综上，本方案以"大车批量推块主攻高分区、小车独立行动补充得分、机械臂专职圆环操作"的三线并行策略为核心，充分发挥子母车架构的并行优势，在有限比赛时间内最大化得分效率，同时兼顾对抗赛的路径封锁与区域占位需求，形成了一套攻守兼备的完整竞赛解决方案。

    \vspace*{\fill}
\end{reportframe}

\end{document}
